\scanpage{230}\mspage{85}
\[
\begin{aligned}
\overset{\smash{\tikz{\draw (0,0)--(0,.08)--(.95,.08)--(.95,0);}}}{\mathbb L_{\deleted{\unsure{w}}W}(A)}
 &=\{(p:Y'\to Y)\in\deleted{\unsure{\operatorname{Fl}}}\operatorname{Fl}(M)
 \mid p^*:\operatorname{Cat}_{/Y}\unsure{-|\to}\operatorname{Cat}_{/Y'}\ \text{transforme}\\
 &\hspace{6em}\text{flèches de }A\text{ en flèches \uuline{de $W$}}\}.
\end{aligned}\tag*{(6.4)}
\]
\diagram{\begin{tikzcd}[ampersand replacement=\&,column sep=1.2em,row sep=1em]
X_1\arrow[d,"f"'] \& X'_1\arrow[l]\arrow[d,"f'"]\\
X_2\arrow[d]\arrow[ur,phantom,"\mathrm{cart}"] \& X'_2\arrow[l]\arrow[d]\\
Y\arrow[ur,phantom,"\mathrm{cart}"] \& Y'\arrow[l,"p"]
\end{tikzcd}}
\margin{$f\in A\Rightarrow f'\in W$ \deleted{\unsure{$f'$}}}
\margin{NB la \unsure{clause peut être|clause précédente} \unsure{lue} en termes de diagr. cartésiens, ce qui \unsure{en est} en fait \unsure{indép. de} l'hypothèse de stabilité de $W$ par isoms.}

Il est clair d'autre part que cette partie de $\operatorname{Fl}(M)$
est stable par changt de base\inserted{, \deleted{et par composition}}, \unsure{donc} $\mathbb L_W(A)\in Q$,
\diagram{\begin{tikzcd}[ampersand replacement=\&,row sep=1em,column sep=small]
X_1\arrow[d,"f"'] \& X'_1\arrow[l,"q_1"']\arrow[d,"f'"]\\
X_2\arrow[d]\arrow[ur,phantom,"\mathrm{cart}"] \& X'_2\arrow[l,"q_2"']\arrow[d]\\
Y_1\arrow[d]\arrow[ur,phantom,"\mathrm{cart}"] \& Y'_1\arrow[l,"p_1"']\arrow[d]\\
Y\arrow[ur,phantom,"\mathrm{cart}"] \& Y'\arrow[l,"p"]
\end{tikzcd}}
et que $(A,\mathbb L_W(A))\in R$. Donc
$\mathbb L_W(A)$ est un plus grand élément
de $\mathcal R(A)$, donc on trouve
\[
\mathcal R=\{(A,L)\in P\unsure{\times|<}Q\mid L\subset\mathbb L_W(A)\}.\tag*{(6.5)}
\]

Inversement, partons de $L\in Q$,
et montrons que l'ouvert $\,{}^{t}\mathcal R(L)$ de $P$ a
un plus grand élément. Pour ceci,
posons
\diagram{\begin{tikzcd}[ampersand replacement=\&,column sep=0.7em,row sep=0.7em]
{} \& {} \& X\arrow[d,"f"'] \& X'\arrow[l,"q"']\arrow[d,"f'"]\\
{\text{(6.6)}\quad\unsure{\mathbb A|A}_W(L)=\big\{(f\unsure{\in|:}X\to Y)\in\inserted{W\subset}\operatorname{Fl}(M)\ \big|} \& {\forall\,\deleted{\unsure{p}}\,:} \& Y\arrow[ur,phantom,"\text{\uuline{cart}}"] \& {Y',}\arrow[l,leftrightarrow,"{\unsure{p}}"]\\
{} \& {} \& {} \& {\text{avec}}\\
{} \& {} \& {} \& {p\in L,\ \deleted{\text{\unsure{le changt}}}}\\
{} \& {} \& {} \& {\deleted{\text{\unsure{induit}}}\inserted{\text{on a}}\,f'\in W\big\}.}
\end{tikzcd}}

(NB Quand $L$ contient les identités, comme c'est
le cas notamment s'il est de la forme $\mathbb L_W(A)$,
-- alors il contient les iso, \unsure{du moins} si $A$ les
contient -- alors $\mathbb A_W(L)$ est l'ens des $f\in\operatorname{Fl}(M)$
satisfaisant à la condition énoncée après $\mid$, sans
avoir à spécifier $f\in W$.) On a $\mathbb A_W(L)\in P$
\deleted{\unsure{par} $A$} et \deleted{il est clair que}\inserted{je vérifie que} $(\mathbb A_W(L),L)\in\mathcal R$,
\inserted{(en utilisant la stabilité de $L$ par composition)}.
\unsure{enfin} $(A,L)\in\mathcal R$ implique $A\subset\mathbb A_W(L)$, \unsure{(ce}
\inserted{qui est clair, en utilisant que $A\subset W$)},
donc $\mathbb A_W(L)$ est un plus grand élément de
$\,{}^{t}\mathcal R(L)$.
% archive sheet number: 230
